\section{Introduction}\label{sec:intro}

Machine-learning network intrusion detectors are mostly evaluated on public benchmarks, where many report macro-F1 close to the maximum. Such scores answer a narrow question: how well a model separates the flows of one dataset after they have been labelled, pooled and shuffled. They say little about an IoT gateway, which has a few CPU cores, sees traffic in time order, builds flows itself and must act on its verdicts. The gap between benchmark and operation is long known~\cite{sommer2010outside} and well documented for security machine learning~\cite{arp2022dos,apruzzese2023sok}; a recent study reports a random forest with benchmark F1 0.97 whose recall fell to 0.0095 in deployment~\cite{tang2026hybrid}.

The sources of this gap are usually studied one at a time: label errors~\cite{engelen2021troubleshooting,liu2022error,flood2024smells}, splits that ignore time~\cite{pendlebury2019tesseract,arp2022dos}, and flow exporters that cut flows differently~\cite{pinto2025hera,vitorino2025revisiting}. Live prototypes rarely check that they classify the flows the model was evaluated on, and automatic blocking is rarely scored by the hosts it blocks wrongly.

We audit one ordinary detector along the whole path: a decision tree over 54 NFStream~\cite{aouini2022nfstream} flow statistics, trained on ten captures of CIC-IoT2023~\cite{neto2023ciciot2023}. We propose no new detection method. Instead, we pass the detector through four steps (labels, splits, live gateway, response) and ask at each step what the previous step hid and which check reveals it. Each step exposed such a problem, and each problem has a cheap check. The lead contributions are:
\begin{itemize}
\item \emph{Evidence-based labels and the measured cost of file-name labels} (Section~\ref{sec:labels}). Background traffic is 91--98.5\% of the flows in the web, brute-force and MITM captures. Trained on file-name labels, the tree reaches a true macro-F1 of 0.396 and a Benign FPR of 71.8\%, against 0.725 and 3.8\% for the deployed tree trained on evidence labels.
\item \emph{A live-fidelity check and the failure modes it exposed} (Section~\ref{sec:live}). We match live and offline flows one by one, account for packets, and measure verdict delay under CPU budgets. This exposed a switch-forwarding pitfall, idle flow meters and an NFStream export-ordering defect, reported as nfstream\#262~\cite{nfstream262}.
\item \emph{Response scored in false blocks} (Section~\ref{sec:response}). Blocking on the first verdict blocked 60 benign host pairs; we report the full trade-off over the number of verdicts required per pair.
\end{itemize}
As supporting evidence, Section~\ref{sec:splits} shows what a random split hides: a change of brute-force target within one capture session (a case study on three host pairs) and connections that span the train/test boundary.
Section~\ref{sec:discussion} states threats to validity and condenses the audit into five checks.
